\documentclass[12pt,a4paper]{article}

% ─── Packages ──────────────────────────────────────────────────────────────────
\usepackage[utf8]{inputenc}
\usepackage[T1]{fontenc}
\usepackage[turkish]{babel}
\usepackage{geometry}
\usepackage{graphicx}
\usepackage{hyperref}
\usepackage{listings}
\usepackage{xcolor}
\usepackage{tabularx}
\usepackage{booktabs}
\usepackage{enumitem}
\usepackage{fancyhdr}
\usepackage{titlesec}
\usepackage{mdframed}
\usepackage{amsmath}
\usepackage{lmodern}

\geometry{margin=2.5cm}

% ─── Colors ────────────────────────────────────────────────────────────────────
\definecolor{fuchsia}{RGB}{217, 70, 239}
\definecolor{codeblue}{RGB}{100, 160, 255}
\definecolor{codegray}{RGB}{40, 40, 40}
\definecolor{lightgray}{RGB}{245, 245, 245}
\definecolor{darkbg}{RGB}{20, 20, 30}

% ─── Code Listing Style ────────────────────────────────────────────────────────
\lstset{
  basicstyle=\ttfamily\small,
  backgroundcolor=\color{lightgray},
  keywordstyle=\color{blue}\bfseries,
  commentstyle=\color{gray}\itshape,
  stringstyle=\color{red!60!black},
  numbers=left,
  numberstyle=\tiny\color{gray},
  numbersep=8pt,
  showstringspaces=false,
  breaklines=true,
  frame=single,
  rulecolor=\color{gray!30},
  captionpos=b,
  tabsize=2,
}

% ─── Header / Footer ───────────────────────────────────────────────────────────
\pagestyle{fancy}
\fancyhf{}
\rhead{\textcolor{fuchsia}{\small CodeAlchemist — Teknik Özellik Belgesi}}
\lhead{\small \leftmark}
\rfoot{\small Sayfa \thepage}

% ─── Title Formatting ──────────────────────────────────────────────────────────
\titleformat{\section}[block]{\large\bfseries\color{fuchsia}}{}{0em}{}[\titlerule]
\titleformat{\subsection}[block]{\normalsize\bfseries\color{blue!70!black}}{}{0em}{}

% ─── Hyperlink Setup ───────────────────────────────────────────────────────────
\hypersetup{
  colorlinks=true,
  linkcolor=fuchsia,
  urlcolor=codeblue,
  citecolor=gray,
  pdftitle={CodeAlchemist Teknik Özellik Belgesi},
  pdfauthor={CodeAlchemist Geliştirme Ekibi},
}

% ───────────────────────────────────────────────────────────────────────────────
\begin{document}

% ─── Title Page ────────────────────────────────────────────────────────────────
\begin{titlepage}
  \centering
  \vspace*{2cm}
  {\Huge\bfseries\color{fuchsia} CodeAlchemist}\\[0.5em]
  {\Large\color{gray} Yapay Zeka Destekli Kodlama Asistanı}\\[0.3em]
  {\large\itshape Teknik Özellik Belgesi}\\[2cm]

  \hrule height 1pt
  \vspace{0.4cm}
  {\large\bfseries Uygulanan Özellikler (7 Modül)}\\
  {\small Onboarding • Mobil UX • Arama/Filtreleme • Geri Bildirim •
  Monaco Editör • Kod Diff/Uygula • Proje/Workspace}
  \vspace{0.4cm}
  \hrule height 1pt

  \vspace{2cm}
  \begin{tabular}{rl}
    \textbf{Proje}  & CodeAlchemist v2.0 \\
    \textbf{Tarih}  & Mart 2026 \\
    \textbf{Ortam}  & React 19 + Flask 3 + SQLAlchemy + Monaco Editor \\
    \textbf{Durum}  & \textcolor{green!60!black}{\bfseries Uygulandı}
  \end{tabular}
\end{titlepage}

% ─── Table of Contents ─────────────────────────────────────────────────────────
\tableofcontents
\newpage

% ═══════════════════════════════════════════════════════════════════════════════
\section{Giriş ve Mimari Genel Bakış}
% ═══════════════════════════════════════════════════════════════════════════════

CodeAlchemist, geliştiricilere yapay zeka destekli kod analizi, hata ayıklama ve üretim ortamına hazır PR oluşturma imkânı sunan modern bir kodlama asistanıdır. Bu belgede, uygulamaya eklenen \textbf{7 temel özellik modülü} teknik düzeyde açıklanmaktadır.

\subsection{Teknoloji Yığını}

\begin{tabular}{lll}
\toprule
\textbf{Katman} & \textbf{Teknoloji} & \textbf{Versiyon} \\
\midrule
Frontend & React & 19.2 \\
Editör & Monaco Editor (\texttt{@monaco-editor/react}) & 4.7 \\
Stil & Tailwind CSS & 3.4 \\
Backend & Python / Flask & 3.x \\
ORM & SQLAlchemy (Flask-SQLAlchemy) & 3.x \\
Auth & Flask-JWT-Extended & 4.x \\
Veritabanı & PostgreSQL / SQLite & — \\
AI & Gemini 2.5 Flash, Claude 3.5, GPT-4o & — \\
\bottomrule
\end{tabular}

\subsection{Klasör Yapısı (İlgili Dosyalar)}

\begin{lstlisting}[language=bash, caption={Proje yapısı - özellikle değiştirilen dosyalar}]
code_alchemist/
├── client/src/
│   ├── components/
│   │   ├── OnboardingTour.jsx        # [YENİ] Hoşgeldin sihirbazı
│   │   ├── MonacoCodeEditor.jsx      # [YENİ] Monaco editör bileşeni
│   │   ├── ProjectManager.jsx        # [YENİ] Proje yöneticisi
│   │   ├── HistoryList.jsx           # [GÜNCELLENDİ] Arama eklendi
│   │   └── ChatInterface.jsx         # [GÜNCELLENDİ] Feedback + Monaco
│   ├── App.jsx                       # [GÜNCELLENDİ] Onboarding + Proje
│   └── index.css                     # [GÜNCELLENDİ] Mobil stiller
└── server/
    ├── models.py                     # [GÜNCELLENDİ] Feedback,Project,File
    └── app.py                        # [GÜNCELLENDİ] 3 yeni API grubu
\end{lstlisting}

\newpage

% ═══════════════════════════════════════════════════════════════════════════════
\section{Özellik 10: İnteraktif Onboarding Turu}
% ═══════════════════════════════════════════════════════════════════════════════

\subsection{Motivasyon}
Yeni kullanıcılar uygulamaya ilk girdiklerinde ne yapacaklarını bilememektedir. Boş ekran ile karşılaşan kullanıcının deneyimini iyileştirmek ve elde tutma oranını artırmak için adım adım bir karşılama sihirbazı gereklidir.

\subsection{Uygulama}

\textbf{Dosya:} \texttt{client/src/components/OnboardingTour.jsx}

Bileşen, 5 adımlı bir modal sihirbaz olarak tasarlanmıştır:

\begin{enumerate}[label=\textbf{Adım \arabic*.}]
  \item \textbf{Hoşgeldin} — Uygulama tanıtımı ve genel bakış
  \item \textbf{Model Seç} — AI model seçiminin nasıl yapılacağı
  \item \textbf{Soru Sor} — Chat arayüzü ve dosya yükleme rehberi
  \item \textbf{Kenar Çubuğu} — Geçmiş ve topluluk erişimi
  \item \textbf{GitHub} — Repo bağlama entegrasyonu
\end{enumerate}

\begin{lstlisting}[language=JavaScript, caption={OnboardingTour - Tamamlanma tespiti}]
const TOUR_KEY = 'codealchemist_onboarding_done';

// Sihirbaz gösterilmeli mi?
export const shouldShowOnboarding = () => !localStorage.getItem(TOUR_KEY);

// Tamamlandığında işaretle
const handleComplete = () => {
  localStorage.setItem(TOUR_KEY, 'true');
  onComplete?.();
};
\end{lstlisting}

\textbf{App.jsx entegrasyonu:}

\begin{lstlisting}[language=JavaScript, caption={App.jsx - Onboarding state ve render}]
const [showOnboarding, setShowOnboarding] = useState(
  () => shouldShowOnboarding()
);

// Kayıt sonrası da göster
const handleAuthSuccess = (data) => {
  // ... auth logic ...
  if (shouldShowOnboarding()) setShowOnboarding(true);
};

// JSX içinde:
{showOnboarding && (
  <OnboardingTour onComplete={() => setShowOnboarding(false)} />
)}
\end{lstlisting}

\subsection{Test Yöntemi}
\begin{lstlisting}[language=bash, caption={Onboarding testini sıfırlama}]
# Tarayıcı konsolunda:
localStorage.removeItem('codealchemist_onboarding_done');
# Ardından sayfayı yenile
\end{lstlisting}

\newpage

% ═══════════════════════════════════════════════════════════════════════════════
\section{Özellik 11: Mobil Deneyim İyileştirmeleri}
% ═══════════════════════════════════════════════════════════════════════════════

\subsection{Motivasyon}
Chat arayüzü masaüstü için optimize edilmiş olmakla birlikte, mobil cihazlarda (özellikle 375px-768px arası) ciddi kullanılabilirlik sorunları bulunmaktaydı: mesaj balonları ekranı taşıyor, klavye açılınca viewport kayıyor, butonlar küçük touch target'a sahip oluyordu.

\subsection{Uygulanan CSS Değişiklikleri}

\textbf{Dosya:} \texttt{client/src/index.css}

\begin{lstlisting}[language=CSS, caption={Mobil chat iyileştirmeleri}]
@media (max-width: 768px) {
  /* Mesaj baloncukları daha geniş */
  .max-w-[85%] { max-width: 93% !important; }
  .max-w-[90%] { max-width: 95% !important; }

  /* Klavye açıkken viewport kaymasını engelle */
  .flex.flex-col.h-full { height: 100dvh; }

  /* Geri bildirim butonları - büyük touch target */
  .text-sm.px-2.py-1.rounded-lg {
    padding: 0.375rem 0.625rem !important;
    font-size: 1rem !important;
  }

  /* Model karşılaştırma - tek sütun */
  .grid.grid-cols-2 {
    grid-template-columns: 1fr !important;
  }
}

/* Onboarding için pulse animasyonu */
@keyframes onboarding-pulse {
  0%, 100% { box-shadow: 0 0 0 0 rgba(217, 70, 239, 0.4); }
  50% { box-shadow: 0 0 0 12px rgba(217, 70, 239, 0); }
}
\end{lstlisting}

\begin{tabular}{lll}
\toprule
\textbf{Sorun} & \textbf{Çözüm} & \textbf{Etkilenen Element} \\
\midrule
Mesaj taşması & max-width artırıldı & .max-w-{[}85\%{]} \\
Viewport kayması & \texttt{100dvh} kullanıldı & Chat container \\
Küçük butonlar & padding artırıldı & Feedback butonları \\
2-sütun grid & 1-sütuna düşürüldü & Model Compare \\
\bottomrule
\end{tabular}

\newpage

% ═══════════════════════════════════════════════════════════════════════════════
\section{Özellik 12: Geçmiş Arama ve Filtreleme}
% ═══════════════════════════════════════════════════════════════════════════════

\subsection{Motivasyon}
Kullanıcılar zamanla çok sayıda konuşma biriktirir. Geçmiş konuşmalar arasında başlık bazlı arama yapılamaması, kullanıcıların eski sorularına dönmesini zorlaştırıyordu.

\subsection{Uygulama}

\textbf{Dosya:} \texttt{client/src/components/HistoryList.jsx}

Arama, tamamen istemci taraflı (client-side) çalışmaktadır. Hiçbir API çağrısı gerekmez; mevcut konuşma listesi anlık olarak filtrelenir.

\begin{lstlisting}[language=JavaScript, caption={Client-side filtreleme mantığı}]
const [searchText, setSearchText] = useState('');

const filtered = (conversations || []).filter(item => {
  if (!searchText.trim()) return true;
  const q = searchText.toLowerCase();
  return (item.title || '').toLowerCase().includes(q);
});
\end{lstlisting}

\textbf{Arama çubuğu JSX:}

\begin{lstlisting}[language=JavaScript, caption={Arama input bileşeni}]
<div className="relative mb-3">
  <span className="absolute left-3 top-1/2 -translate-y-1/2 text-xs">
    🔍
  </span>
  <input
    type="text"
    value={searchText}
    onChange={e => setSearchText(e.target.value)}
    placeholder="Search conversations..."
    className="w-full bg-gray-900/60 border border-gray-700 rounded-lg
               pl-8 pr-8 py-2 text-xs text-gray-200 ..."
  />
  {searchText && (
    <button onClick={() => setSearchText('')}>✕</button>
  )}
</div>
\end{lstlisting}

\subsection{Özellikler}
\begin{itemize}
  \item Anlık filtreleme — her tuş vuruşunda güncellenen liste
  \item "X" butonu ile hızlı temizleme
  \item Sonuç bulunamazsa: \textit{No conversations matching "..."} mesajı
  \item Arama yokken tüm liste görünür (sıfır gecikme)
\end{itemize}

\newpage

% ═══════════════════════════════════════════════════════════════════════════════
\section{Özellik 14: AI Yanıt Kalitesi Geri Bildirimi}
% ═══════════════════════════════════════════════════════════════════════════════

\subsection{Motivasyon}
Kullanıcıların AI yanıtlarını değerlendirebilmesi hem UX kalitesini ölçmek hem de ileride model kalibrasyonu için veri toplamak açısından kritik öneme sahiptir.

\subsection{Veritabanı Modeli}

\textbf{Dosya:} \texttt{server/models.py}

\begin{lstlisting}[language=Python, caption={Feedback veritabanı modeli}]
class Feedback(db.Model):
    """AI yanıtları için kullanıcı geri bildirimi (👍/👎)"""
    id         = db.Column(db.Integer, primary_key=True)
    history_id = db.Column(db.Integer, db.ForeignKey('history.id'))
    user_id    = db.Column(db.Integer, db.ForeignKey('user.id'), nullable=True)
    rating     = db.Column(db.Integer, nullable=False)  # +1 veya -1
    created_at = db.Column(db.DateTime, default=datetime.utcnow)

    __table_args__ = (
        db.UniqueConstraint('history_id', 'user_id',
                           name='_feedback_unique_uc'),
    )
\end{lstlisting}

\subsection{API Endpoint'leri}

\textbf{Dosya:} \texttt{server/app.py}

\begin{tabular}{lll}
\toprule
\textbf{Metod} & \textbf{Yol} & \textbf{Açıklama} \\
\midrule
POST & \texttt{/api/feedback} & Oy gönder / toggle / güncelle \\
GET  & \texttt{/api/feedback/<id>} & 👍/👎 sayısı ve kullanıcı oyu \\
\bottomrule
\end{tabular}

\textbf{Toggle mantığı:}

\begin{lstlisting}[language=Python, caption={POST /api/feedback - toggle ve güncelleme}]
existing = Feedback.query.filter_by(
    history_id=history_id, user_id=user_id).first()
if existing:
    if existing.rating == rating:
        db.session.delete(existing)   # Aynı oy → geri al
    else:
        existing.rating = rating       # Farklı oy → güncelle
else:
    db.session.add(Feedback(...))      # Yeni oy
\end{lstlisting}

\subsection{Frontend Uygulama}

\textbf{Dosya:} \texttt{client/src/components/ChatInterface.jsx}

\begin{lstlisting}[language=JavaScript, caption={Optimistic update ile feedback butonu}]
const handleFeedback = async (historyId, rating) => {
  const current = feedbacks[historyId];
  const newRating = current === rating ? null : rating; // toggle
  setFeedbacks(prev => ({ ...prev, [historyId]: newRating }));

  try {
    await fetch(`${apiBase}/api/feedback`, {
      method: 'POST',
      headers: { 'Content-Type': 'application/json', ...authHeaders },
      body: JSON.stringify({ history_id: historyId, rating })
    });
  } catch (err) {
    setFeedbacks(prev => ({ ...prev, [historyId]: current })); // geri al
  }
};
\end{lstlisting}

\newpage

% ═══════════════════════════════════════════════════════════════════════════════
\section{Özellik 15: Monaco Kod Editörü Entegrasyonu}
% ═══════════════════════════════════════════════════════════════════════════════

\subsection{Motivasyon}
Kullanıcı kodu yalnızca düz bir \texttt{<textarea>} alanına yapıştırabiliyordu. VS Code'un motorunu kullanan Monaco Editor entegrasyonu şunları sunar: sözdizimi vurgulaması, otomatik tamamlama, satır numaraları, katlanabilir bloklar ve çoklu dil desteği.

\subsection{Bağımlılık}

\begin{lstlisting}[language=bash, caption={Zaten yüklü olan paket}]
"@monaco-editor/react": "^4.7.0"
\end{lstlisting}

\subsection{MonacoCodeEditor Bileşeni}

\textbf{Dosya:} \texttt{client/src/components/MonacoCodeEditor.jsx}

\begin{lstlisting}[language=JavaScript, caption={Monaco bileşeni temel yapısı}]
import Editor from '@monaco-editor/react';

const MonacoCodeEditor = ({
  value, onChange, language = 'python',
  onLanguageChange, theme = 'dark', height = '200px'
}) => {
  return (
    <Editor
      height={height}
      language={language}
      value={value}
      onChange={(val) => onChange?.(val ?? '')}
      theme={theme === 'dark' ? 'vs-dark' : 'light'}
      options={{
        minimap: { enabled: false },
        fontSize: 13,
        fontFamily: "'JetBrains Mono', Consolas, monospace",
        fontLigatures: true,
        wordWrap: 'on',
        quickSuggestions: true,
        acceptSuggestionOnEnter: 'smart',
      }}
    />
  );
};
\end{lstlisting}

\subsection{ChatInterface.jsx Entegrasyonu}
Kullanıcı \textbf{💻 Code} butonuna tıklayarak Monaco editörüne, \textbf{⌨️ Plain} butonuna tıklayarak düz textarea'ya geçiş yapabilir:

\begin{lstlisting}[language=JavaScript, caption={Monaco/Textarea toggle}]
const [useMonacoEditor, setUseMonacoEditor] = useState(false);

{useMonacoEditor ? (
  <MonacoCodeEditor
    value={question}
    onChange={setQuestion}
    language={monacoLanguage}
    onLanguageChange={setMonacoLanguage}
    theme="dark"
    height="180px"
  />
) : (
  <textarea value={question} onChange={...} />
)}
\end{lstlisting}

\subsection{Desteklenen Diller}
Python, JavaScript, TypeScript, Go, Rust, Java, C/C++, SQL, HTML, CSS, JSON, Markdown ve Düz Metin.

\newpage

% ═══════════════════════════════════════════════════════════════════════════════
\section{Özellik 16: Kod Diff ve Uygula Sistemi}
% ═══════════════════════════════════════════════════════════════════════════════

\subsection{Motivasyon}
AI bir kod düzeltmesi önerdiğinde, kullanıcı değişikliği elle kopyalamak zorundaydı. \texttt{react-diff-viewer-continued} kütüphanesi zaten bulunduğundan, bu kütüphane üzerine inşa edilen "Apply" butonu iş akışını dramatik biçimde hızlandırır.

\subsection{Mevcut Altyapı}
\texttt{ChatInterface.jsx}, AI yanıtlarındaki kod bloklarını \texttt{SmartMarkdown} bileşeni üzerinden işlemektedir. Her kod bloğu zaten sözdizimi vurgulu bir görünüme sahiptir. Diff viewer \texttt{react-diff-viewer-continued@4.1.2} ile desteklenmektedir.

\subsection{Apply Patch İş Akışı}

\begin{lstlisting}[language=JavaScript, caption={Kod bloğundan patch çıkarma}]
// AI yanıtındaki ```...``` bloklarını çıkarır
const extractCodeBlocks = (responseText) => {
  const regex = /```(\w+)?\n([\s\S]*?)```/g;
  const blocks = [];
  let match;
  while ((match = regex.exec(responseText)) !== null) {
    blocks.push({ language: match[1] || 'text', content: match[2] });
  }
  return blocks;
};
\end{lstlisting}

\begin{lstlisting}[language=JavaScript, caption={Diff görüntüleme bileşeni}]
import ReactDiffViewer from 'react-diff-viewer-continued';

<ReactDiffViewer
  oldValue={originalCode}      // Kullanıcının orijinal kodu
  newValue={suggestedCode}     // AI'nın önerdiği kod
  splitView={true}
  useDarkTheme={true}
  showDiffOnly={false}
/>
\end{lstlisting}

\textbf{Apply Patch} butonuna tıklandığında:
\begin{itemize}
  \item Kullanıcının kod alanı (Monaco veya textarea) güncellenir
  \item Değişiklik görsel olarak onaylanır (toast bildirimi)
  \item Chat bağlamı AI'ya iletilerek ardışık sorularla devam edilebilir
\end{itemize}

\newpage

% ═══════════════════════════════════════════════════════════════════════════════
\section{Özellik 17: Proje / Workspace Kavramı}
% ═══════════════════════════════════════════════════════════════════════════════

\subsection{Motivasyon}
Kullanıcılar tek bir konuşmada izole şekilde çalışabiliyordu. Gerçek dünya geliştirme senaryolarında birden fazla dosyanın (ör. \texttt{App.jsx} + \texttt{app.py}) aynı bağlamda sunulması gerekir.

\subsection{Veritabanı Modelleri}

\textbf{Dosya:} \texttt{server/models.py}

\begin{lstlisting}[language=Python, caption={Project ve ProjectFile modelleri}]
class Project(db.Model):
    id          = db.Column(db.Integer, primary_key=True)
    user_id     = db.Column(db.Integer, db.ForeignKey('user.id'))
    name        = db.Column(db.String(255), nullable=False)
    description = db.Column(db.Text, nullable=True)
    created_at  = db.Column(db.DateTime, default=datetime.utcnow)
    updated_at  = db.Column(db.DateTime, onupdate=datetime.utcnow)

    files = db.relationship('ProjectFile', backref='project',
                            lazy='dynamic', cascade='all, delete-orphan')


class ProjectFile(db.Model):
    id         = db.Column(db.Integer, primary_key=True)
    project_id = db.Column(db.Integer, db.ForeignKey('project.id'))
    name       = db.Column(db.String(255), nullable=False)
    content    = db.Column(db.Text, nullable=False, default='')
    language   = db.Column(db.String(50), default='plaintext')
    created_at = db.Column(db.DateTime, default=datetime.utcnow)
\end{lstlisting}

\subsection{API Endpoint'leri}

\textbf{Dosya:} \texttt{server/app.py}

\begin{tabular}{llp{7cm}}
\toprule
\textbf{Metod} & \textbf{Yol} & \textbf{Açıklama} \\
\midrule
GET    & \texttt{/api/projects}                   & Projeleri listele \\
POST   & \texttt{/api/projects}                   & Yeni proje oluştur \\
DELETE & \texttt{/api/projects/<id>}              & Projeyi sil \\
GET    & \texttt{/api/projects/<id>/files}        & Dosya listesi \\
POST   & \texttt{/api/projects/<id>/files}        & Dosya ekle \\
DELETE & \texttt{/api/projects/<id>/files/<fid>} & Dosyayı sil \\
GET    & \texttt{/api/projects/<id>/context}      & AI bağlamı (tüm dosyalar birleşik) \\
\bottomrule
\end{tabular}

\subsection{AI Bağlamı Oluşturma}

\begin{lstlisting}[language=Python, caption={/api/projects/<id>/context - Birleşik AI bağlamı}]
@app.route('/api/projects/<int:project_id>/context')
@jwt_required()
def get_project_context(project_id):
    project = Project.query.filter_by(id=project_id, ...).first_or_404()
    files = project.files.order_by(ProjectFile.name).all()

    context_parts = [f"# Project: {project.name}"]
    for f in files:
        context_parts.append(
            f"\n## File: {f.name} ({f.language})\n```{f.language}\n{f.content}\n```"
        )
    return jsonify({'context': '\n'.join(context_parts)})
\end{lstlisting}

\subsection{Frontend: ProjectManager Bileşeni}

\textbf{Dosya:} \texttt{client/src/components/ProjectManager.jsx}

\begin{lstlisting}[language=JavaScript, caption={ProjectManager - temel işlevler}]
// Proje seçildiğinde ana App bileşenine bildirim
const handleSelectAndClose = (project) => {
  onSelectProject?.(project);     // App.jsx'te activeProject güncellenir
  onClose?.();
};

// Seçili projenin tüm dosyaları chat context'e enjekte edilir
const context = await fetch(`/api/projects/${id}/context`);
\end{lstlisting}

\textbf{App.jsx entegrasyonu:}

\begin{lstlisting}[language=JavaScript, caption={App.jsx - 📁 Projects sekmesi}]
{showProjectManager && (
  <ProjectManager
    apiBase={API_BASE}
    authHeaders={token ? { Authorization: `Bearer ${token}` } : {}}
    onSelectProject={(project) => setActiveProject(project)}
    activeProjectId={activeProject?.id}
    onClose={() => setShowProjectManager(false)}
  />
)}
\end{lstlisting}

\newpage

% ═══════════════════════════════════════════════════════════════════════════════
\section{Özet ve Performans Değerlendirmesi}
% ═══════════════════════════════════════════════════════════════════════════════

\subsection{Uygulanan Özellikler Özeti}

\begin{tabularx}{\textwidth}{lXll}
\toprule
\textbf{\#} & \textbf{Özellik} & \textbf{Frontend} & \textbf{Backend} \\
\midrule
10 & Onboarding Turu       & \textcolor{green!60!black}{✓} & —  \\
11 & Mobil UX              & \textcolor{green!60!black}{✓} & —  \\
12 & Arama \& Filtreleme   & \textcolor{green!60!black}{✓} & —  \\
14 & Geri Bildirim 👍/👎   & \textcolor{green!60!black}{✓} & \textcolor{green!60!black}{✓} \\
15 & Monaco Editör         & \textcolor{green!60!black}{✓} & —  \\
16 & Kod Diff \& Apply     & \textcolor{green!60!black}{✓} & —  \\
17 & Proje/Workspace       & \textcolor{green!60!black}{✓} & \textcolor{green!60!black}{✓} \\
\bottomrule
\end{tabularx}

\subsection{Yeni Veritabanı Tabloları}

\begin{tabular}{lll}
\toprule
\textbf{Tablo} & \textbf{Amaç} & \textbf{Önemli Kısıt} \\
\midrule
\texttt{feedback}     & 👍/👎 oyları      & UNIQUE (history\_id, user\_id) \\
\texttt{project}      & Workspace meta    & user\_id FK \\
\texttt{project\_file} & Dosya içerikleri & project\_id FK, cascade delete \\
\bottomrule
\end{tabular}

\subsection{Yeni API Endpoint Sayısı}

\begin{itemize}
  \item Feedback API: \textbf{2 endpoint} (POST, GET)
  \item Project API:  \textbf{7 endpoint} (CRUD + context)
  \item Toplam yeni endpoint: \textbf{9}
\end{itemize}

\subsection{Bundle Etkisi}

Monaco Editor, lazy load ile yüklendiğinden başlangıç bundle boyutunu artırmaz. \texttt{@monaco-editor/react}, editör DOM'a mount olduğunda dinamik olarak yüklenir. Tahmini ek yük: ilk Monaco açılışında $\approx$2MB.

% ═══════════════════════════════════════════════════════════════════════════════
\section{Sonuç}
% ═══════════════════════════════════════════════════════════════════════════════

Bu belgede CodeAlchemist uygulamasına eklenen 7 özellik modülü teknik olarak açıklanmıştır. Tüm değişiklikler mevcut mimari ile uyumlu biçimde geliştirilmiş, veritabanı şeması geri çekilebilir (backward-compatible) şekilde genişletilmiş ve frontend bileşenleri render perfomansına dikkat edilerek tasarlanmıştır.

\begin{mdframed}[backgroundcolor=lightgray]
\textbf{Önemli Not:} Tüm yeni özellikler, mevcut bir Flask uygulamasına minimal müdahaleyle eklenmiştir. Geri bildirim sistemi bağımsız tabloya yazıldığından, gelecekte makine öğrenimi modeli ince ayarı (fine-tuning) için hazır veri kümesi sunar.
\end{mdframed}

\vspace{1cm}
\begin{center}
\textcolor{fuchsia}{\rule{10cm}{0.5pt}}\\[0.3em]
\textit{CodeAlchemist — Yapay Zeka Destekli Kodlama Asistanı}\\
\textit{Teknik Belge v2.0 — Mart 2026}
\end{center}

\end{document}
